\section{Parallel implementation}
\label{sec:parallel}

\Cref{prop:order} says the search tree may be explored in any order, which
leaves the traversal free to be organised around the processor rather than
around the query. Both implementations use that freedom, and they use it in
opposite directions. The host kernel fills the lanes of one vector register with
boxes drawn from different queries, so a lane never idles while its neighbours
work (\cref{sec:parallel:host}). The device kernel does the reverse: it moves
boxes \emph{out} of the block that produced them, into a queue that the next
launch spreads over the whole device (\cref{sec:parallel:device}). The broad
phase needs no such treatment, since its parallelism is already the outermost
loop of \cref{alg:cell,alg:cellmin} and its count-then-fill output carries no
contention.

What both designs answer is a single property of the problem: the cost of a CCD
query is heavy-tailed. Most candidate pairs are settled by the first box the
search looks at, while a few require tens of thousands or more
(\cref{tab:difficulty}, at the end of this section). A design that pins one
query to one execution resource --- a lane, a thread, a block --- is therefore
load-imbalanced by construction, and not by a factor of two. The remainder of
this section presents the two kernels and the parameters they were measured at;
\cref{sec:parallel:why} returns to the measurement and to what the two
distributions say about where the work goes.

\subsection{Host: lanes packed across queries}
\label{sec:parallel:host}

The unit of work on the host is the \emph{box}, not the query. A vector of $V$
lanes ($V=8$ on AVX-512 and AVX2, $4$ on NEON) holds $V$ boxes that may belong
to $V$ different queries, and the kernel advances all of them through one
iteration of \cref{alg:np} together. Queries are processed in blocks of $V$
consecutive candidates, one block per task of the parallel loop; a block's
geometry, per-axis tolerances and certified error bound are computed once on
entry, and each query contributes the root box $[0,\toi]\times[0,1]^{2}$ to a
work list private to the block. A lane that decides its box takes the next box
off that list, whichever query it belongs to, and the lane's geometry is
reloaded only when the query actually changes. \Cref{prop:order} is what makes
this legal: the answer does not depend on the order in which boxes are visited,
so interleaving unrelated queries in one register cannot change it.

One iteration proceeds in three phases, and the division is what allows a single
flat vector sweep to serve lanes that are doing different things.

\emph{Classify} decides the fate of every lane from quantities it already
carries. Each box stores the value of $\Fdiff$ at its eight corners, so the
componentwise hull of \cref{eq:hull} is a min and a max over stored values, taken
across all lanes at once, and no function evaluation occurs in this phase at all.
A lane whose hull excludes the origin by more than the certified bound is
rejected and goes idle; a lane that the acceptance test of
\cref{eq:accept-tight} accepts, or that has reached the depth cap, or whose box
can no longer be bisected in floating point, lowers its query's running minimum
to the box's $t$ lower bound and goes idle. Every remaining lane chooses its
split axis, the one furthest past its own tolerance, and the midpoint of that
axis.

\emph{Evaluate} is the only phase that computes $\Fdiff$. A bisection cuts the
parent at a mid-face, so the two children share that face and each inherits its
outer face from the parent: the four corner values of the mid-face are all that
the split needs. Those four parameter triples are laid out for every splitting
lane and evaluated in one vector sweep, even though the lanes are splitting
different axes --- the axis enters only through which parent coordinates are
written into the sweep's input, so the arithmetic is identical across lanes.
Lanes that are not splitting ride along and their results are discarded, which
is cheaper than branching inside the loop.

\emph{Assemble} forms the two children from the parent's corners and the newly
evaluated mid-face. The upper child goes to the work list; the lower child stays
in the lane, overwriting its parent in place. Keeping the lower child is what
makes each lane depth-first in $t$: the child that may contain an earlier impact
is examined first, so the running minimum falls sooner and prunes harder, the
work list carries half the traffic, and the lane needs no refill and no geometry
reload. Either child is dropped before it is written if the running minimum has
already passed its $t$ lower bound, or --- for a vertex-face query, whose domain
is the triangle $u+v\le1$ rather than the unit square --- if the child lies
outside that domain. When the block is exhausted, its per-query minima are either
written out per pair or folded into the shared minimum by compare-and-swap,
according to what the caller asked for.

\begin{algorithm}[!tb]
\caption{Host narrow phase. Lanes hold boxes from different queries, which
\cref{prop:order} permits. $V$ is the vector width and $D$ the depth cap.}
\label{alg:host}
\footnotesize
\begin{algorithmic}[1]
\Function{HostNarrowPhase}{candidates $\Cand$}
  \ForAll{blocks $Q$ of $V$ consecutive queries \textbf{in parallel}}
    \State load geometry, $\tau$ and $\epsilon$ per query;\;
           $\toi_{q} \gets$ shared minimum, or $1$ when the caller wants one
           answer per pair
    \State $S \gets$ work list holding each query's root box with its $8$
           corner values
    \State all $V$ lanes idle
    \While{true}
      \ForAll{idle lanes $l$}
        \Comment{refill}
        \State pop $\dom$ from $S$, discarding any with
               $\dom.t_{\ell} \ge \toi_{\dom.q}$
        \State bind $\dom$ to lane $l$, reloading lane geometry only if
               $\dom.q$ changed
      \EndFor
      \State \textbf{if} no lane holds a box \textbf{then break}
      \State $[\underline{F},\overline{F}] \gets$ hull of the $8$ stored
             corners, per component, all lanes at once
      \ForAll{lanes $l$ holding a box}
        \Comment{phase 1: classify, evaluating nothing}
        \If{$\dom_{l}$ is rejected by the test of \cref{eq:err}}
          \State lane $l \gets$ idle
        \ElsIf{\Call{Accept}{$\dom_{l}$} \textbf{or} $\mathrm{depth}_{l} \ge D$
               \textbf{or} $\dom_{l}$ cannot be bisected}
          \State $\toi_{q_{l}} \gets \min(\toi_{q_{l}},\ \dom_{l}.t_{\ell})$;\;
                 lane $l \gets$ idle
        \Else
          \State $a_{l} \gets \arg\max_{k} \dom_{l}.w_{k}/\tau_{k}$;\;
                 $m_{l} \gets$ midpoint of axis $a_{l}$
        \EndIf
      \EndFor
      \State lay out the $4$ mid-face corner parameters of every splitting lane
      \State evaluate $\Fdiff$ there in \emph{one} vector sweep across all lanes
        \Comment{phase 2}
      \ForAll{splitting lanes $l$}
        \Comment{phase 3: assemble}
        \State push the upper child onto $S$: the parent's outer face, the new
               mid face
        \State keep the lower child in lane $l$, overwritten in place
          \Comment{depth-first in $t$}
        \State discard either child whose $t$ lower bound has passed
               $\toi_{q_{l}}$, or which leaves the query's domain
      \EndFor
    \EndWhile
    \State write $\toi_{q}$ out, or lower the shared minimum by
           compare-and-swap
  \EndFor
\EndFunction
\end{algorithmic}
\end{algorithm}

\subsection{Device: a small stack, and a queue that redistributes}
\label{sec:parallel:device}

The device kernel runs one thread per query and one launch per round of work. A
launch is given a set of boxes, searches them to completion without ever waiting
for another launch, and writes the boxes it could not keep into a queue in global
memory; the host then relaunches the same kernel on that queue. Within a launch
the unit of work is again the box: a thread that finishes its own query's subtree
takes a box belonging to any query from a stack its block shares, and reloads
geometry only if the query changed. \Cref{fig:gpu-workflow} shows the two levels
of storage and \cref{alg:device} gives both the kernel and the host loop around
it.

The per-block stack lives in shared memory and holds $64$ entries, which is
small on purpose. The stack is where a query's subtree accumulates, so its
capacity decides how long a block keeps a heavy query to itself: with room for a
thousand boxes a single block grinds through a query of tens of thousands
(\cref{tab:difficulty}) while its neighbours finish and idle, and no other block
can help, because nothing the block holds is visible outside it. Spilling after $64$ entries pushes that
subtree into the global queue instead, where the next launch distributes it over
every block on the device. The effect is large and it is a scheduling effect, not
a pruning one: the false positives and the reported times of impact are identical
at every capacity (\cref{sec:parallel:tuning}).

The global queue is double buffered. A launch reads one half and writes the
other, never both, so a producer and its consumer are always in different
launches and no thread ever waits on another thread's progress. This is the
reason the queue is organised in rounds at all rather than as a single shared
work pool: a pool whose writers spin on a compare-and-swap until a reader frees a
slot deadlocks as soon as the queue is driven hard, which is precisely the
condition that redistributing work creates. The host loop therefore alternates:
it reads the write cursor after each launch and relaunches on what that launch
produced, swapping the roles of the two halves, until a round produces nothing.

A box that does not fit in the write buffer is counted and discarded. Discarding
is sound rather than merely convenient: every exit from the search other than a
rejection accepts the box at its $t$ lower bound, so a box never examined can
only leave the reported time of impact earlier than the truth, never later
(\cref{sec:np}) --- the failure direction that matters is closed by
construction, and what an undersized queue costs is accuracy. The host reads the
deficit, enlarges the queue by exactly that many entries and reruns the batch
from its seeds, keeping the times of impact already found so that the rerun
prunes harder. Growth is by the deficit rather than geometric because every call
clears the queue it is given, so an oversized queue is paid for on every
subsequent call and not once; the retry is bounded at $32$ rounds, after which
the kernel reports the shortfall and stops, since giving up is conservative and
spinning is not finite.

A work list that is meant never to overflow can be sized from the search rather
than from a measurement, which is what the vertex-quad kernel does with its
per-thread stack: a pop pushes at most $k$ children, so a depth-$D$ search needs
at most $1+kD$ entries for one query, and the headroom is nearly free --- a
$257$-entry frame is $14.7\,\mathrm{KB}$ of which only the first few entries are
ever touched, and register count and spill are unchanged between capacities $4$
and $513$. The shared stack above is the opposite case and is sized against that
rule on purpose: it is not meant to hold a query's frontier, only to pass it on.

The device broad phase is a more conventional port of \cref{alg:cell,alg:cellmin}:
count, prefix sum and scatter, with a radix sort and a device scan from
CUB~\citep{cub} in place of their host equivalents, one box per thread, and a
scalar inner loop where the host uses its $32$-lane mask. Its shape suits a GPU
because no part of it walks a window sequentially.

\begin{figure}[htbp]
  \centering
  \resizebox{0.86\linewidth}{!}{% The device narrow phase's work flow: a small per-block shared stack that spills
% early into a double-buffered global queue, which the next launch redistributes.
%
% Every node carries an explicit inner sep so its text keeps a margin from its
% frame, group labels are placed outside the group they name, and arrows are
% shortened at both ends so no head or tail lands on a border.
\begin{tikzpicture}[
  font=\footnotesize,
  >={Stealth[length=2mm]},
  every edge/.append style={shorten >=1.5pt, shorten <=1.5pt},
  blk/.style={draw=sccdInk, rounded corners=1pt, minimum width=19mm,
              minimum height=7mm, align=center, fill=white,
              inner xsep=6pt, inner ysep=4pt},
  stk/.style={draw=sccdTeal, fill=sccdTeal!8, rounded corners=1pt,
              minimum width=19mm, minimum height=8mm, align=center,
              inner xsep=6pt, inner ysep=4pt},
  buf/.style={draw=sccdPlum, fill=sccdPlum!8, rounded corners=1pt,
              minimum width=32mm, minimum height=9mm, align=center,
              inner xsep=7pt, inner ysep=5pt},
  lbl/.style={font=\scriptsize\itshape, text=sccdInk, align=left, inner sep=1pt},
  grp/.style={draw=sccdGrid, rounded corners=2pt, inner sep=4.5mm},
  flow/.style={->, shorten >=1.5pt, shorten <=1.5pt},
]

% ======================= launch k =======================================
\foreach \i/\x in {0/0, 1/2.4, 2/4.8} {
  \node[blk] (b\i) at (\x,1.9) {block \i};
  \node[stk] (s\i) at (\x,0.85) {shared stack\\[1pt]\scriptsize 64 entries};
  \draw[flow, draw=sccdInk] (b\i) -- (s\i);
}
\node[grp, fit=(b0)(s2)] (g1) {};
\node[lbl, anchor=south west] at ([yshift=1.2mm]g1.north west) {launch $k$};

% The pathological query arrives from outside the group.
\node[lbl, text=sccdAmber, align=center, anchor=south]
  (heavy) at ([yshift=1.2mm]g1.north) {a query needing $>10^{6}$ boxes};
\draw[flow, draw=sccdAmber, thick] (heavy.south) -- (b1.north);

% ======================= spill ==========================================
\node[buf] (wbuf) at (2.4,-0.75)
  {global queue, buffer $A$\\[1pt]\scriptsize launch $k$ \emph{writes} it};
\foreach \i in {0,1,2} {
  \draw[flow, draw=sccdPlum, dashed] (s\i) -- (wbuf);
}
\node[lbl, text=sccdPlum, anchor=west] at ([xshift=5mm]wbuf.east)
  {spill when full;\\overflow is counted,\\not written};

% ======================= swap ===========================================
\node[buf, draw=sccdMoss, fill=sccdMoss!8] (rbuf) at (2.4,-2.45)
  {the same buffer $A$\\[1pt]\scriptsize launch $k{+}1$ \emph{reads} it};
\draw[flow, draw=sccdInk, thick] (wbuf) -- (rbuf);
\node[lbl, anchor=west] at ([xshift=2mm]$(wbuf.south)!0.5!(rbuf.north)$)
  {launch $k$ ends; the buffers swap roles};

% ======================= launch k+1 =====================================
\foreach \i/\x in {0/0, 1/2.4, 2/4.8} {
  \node[blk] (c\i) at (\x,-4.3) {block \i};
  \draw[flow, draw=sccdMoss] (rbuf) -- (c\i);
}
\node[grp, fit=(c0)(c2)] (g2) {};
\node[lbl, anchor=north west] at ([yshift=-1.2mm]g2.south west)
  {launch $k{+}1$};
\node[lbl, text=sccdMoss, anchor=west] at ([xshift=5mm]g2.east)
  {the heavy query's subtree\\is now spread over\\every block};
\node[lbl, anchor=west] at ([xshift=5mm]rbuf.east)
  {and writes buffer $B$,\\which launch $k{+}2$ reads:\\the two never meet};

\begin{scope}[on background layer]
  \node[grp, fit=(b0)(s2)] {};
  \node[grp, fit=(c0)(c2)] {};
\end{scope}
\end{tikzpicture}
}
  \caption{Where a box lives in the device narrow phase. Each block keeps a
    $64$-entry stack in shared memory, visible only to its own threads, and
    spills beyond that into the half of the global queue that the current launch
    writes. The launch that follows reads that half and writes the other, so a
    producer and its consumer are never in the same launch and nothing spins. A
    subtree that crosses into the queue becomes available to every block, where
    a larger shared stack would have kept it inside the block that seeded it.}
  \label{fig:gpu-workflow}
\end{figure}

\begin{algorithm}[!tb]
\caption{Device narrow phase: a kernel that runs to completion on the work it
was given, and a host loop that relaunches it on what it spilled. $S$ is the
per-block shared stack ($64$ entries) and $G$ the double-buffered global queue.}
\label{alg:device}
\footnotesize
\begin{algorithmic}[1]
\Function{Kernel}{read buffer $G_{\mathrm{in}}$, write buffer $G_{\mathrm{out}}$}
  \State thread $t$ takes one box: its own query's root, or one entry of
         $G_{\mathrm{in}}$
  \While{any thread of the block holds a box \textbf{or} $S \ne \emptyset$}
    \State refresh the block's bound from the global minimum
      \Comment{every iteration; \cref{sec:parallel:tuning}}
    \State \textbf{if} $\dom.t_{\ell} \ge$ bound \textbf{then} idle this thread
    \If{$\mathrm{depth} \ge D$ \textbf{or} $\dom$ cannot be bisected}
      \State lower $\toi$ to $\dom.t_{\ell}$; idle this thread
        \Comment{exhaustion accepts}
    \Else
      \State $\dom.t_{h} \gets \min(\dom.t_{h},\ \mathrm{bound})$
      \State $(\dom^{-},\dom^{+}) \gets$ \Call{Split}{$\dom$}
      \State evaluate \textbf{all $8$ corners of each child}
        \Comment{all eight; \cref{alg:host} inherits four}
      \State discard a child the bound has passed or which leaves the domain
      \State \textbf{if} \Call{Accept}{} holds for a child \textbf{then} lower
             $\toi$ to its $t_{\ell}$ and discard it
      \State \textbf{if} both children remain \textbf{then} descend into the
             smaller $t_{\ell}$ and set the other aside
      \State \textbf{else if} one remains \textbf{then} descend into it
             \textbf{else} idle this thread
    \EndIf
    \If{a box was set aside}
      \State claim a slot in $S$; if $S$ is full, bump-allocate in
             $G_{\mathrm{out}}$
      \State if $G_{\mathrm{out}}$ is full, add to the deficit and
             \textbf{discard}
        \Comment{reports early, never late}
    \EndIf
    \State idle threads claim from $S$ by compare-and-swap, reloading geometry
           on a query change
  \EndWhile
\EndFunction
\Statex
\Function{Host}{candidates $\Cand$}
  \Repeat
    \State \Call{Kernel}{seeds, $A$}
      \Comment{one thread per query}
    \While{the buffer just written is non-empty}
      \State \Call{Kernel}{it, the other one}; swap the two
        \Comment{reader and writer never meet}
    \EndWhile
  \Until{the deficit is zero \textbf{or} $32$ rounds have been spent}
  \State between rounds: grow $G$ by the deficit, keeping the times of impact
         found so far
\EndFunction
\end{algorithmic}
\end{algorithm}

\subsection{Precision}
\label{sec:parallel:precision}

Both narrow phases compute in double precision internally whatever type the
caller stores geometry in, and this is a requirement of the method: in single
precision the certified bound of \cref{eq:err} and the tolerances of
\cref{eq:tol} sit too close together for the search to keep its guarantee, and
the single-precision device kernel reported nine times of impact \emph{later}
than the true root on armadillo-rollers, violating \cref{def:conservative}.

The interface stays templated on the storage type, so a caller with float
geometry keeps float buffers. Widening on load is exact, whereas narrowing the
computed time of impact back to float is not: round-to-nearest can round
\emph{up}, publishing a value later than the one the search proved, so the
narrowing rounds toward $-\infty$. Reporting earlier is always safe; reporting
later is the failure the whole design exists to prevent.

\subsection{Parameters settled by measurement}
\label{sec:parallel:tuning}

Five quantities above are constants chosen by measurement rather than by
derivation. They are collected here so that the two kernels can be read without
them, and reported because a reimplementation has to choose them too.

\emph{Vector width.} The two query types peak at different widths, since a
vertex-face corner costs one more multiply-add and one more subtraction than an
edge-edge corner and its sweeps therefore take longer to amortise the per-lane
scalar work between them. On NEON, over cloth-funnel, narrowing vertex-face from
four lanes to two moves it from $297.7$ to $246.0\,\mathrm{ms}$, and widening
edge-edge from two to four moves it from $599.9$ to $494.6\,\mathrm{ms}$; eight
and sixteen lanes are worse for both. The x86 widths quoted in
\cref{sec:parallel:host} follow from the same operation count and have not been
measured. One caveat generalises beyond this kernel: the two types must be timed
in separate processes, because timed together the edge-edge figure varies by
about $20\%$ with the lane width chosen for the \emph{vertex-face}
instantiation, which the measurement never executes. That is an effect of code
layout rather than of work, and it is large enough to manufacture or to conceal
a difference of the size reported here.

\emph{Shared-stack capacity.} Reducing the per-block capacity from $1024$
entries to $256$ and then to $64$ moves the device narrow phase from $4059.2$ to
$198.2$ and then to $37.2\,\mathrm{ms}$ on cloth-funnel, from $3133.7$ to
$277.2$ and $75.8\,\mathrm{ms}$ on armadillo-rollers, and from $129.8$ to
$110.2$ and $105.9\,\mathrm{ms}$ on cloth-ball. False positives and reported
times of impact are identical at every capacity, which is what identifies the
effect as redistribution and not pruning. The smaller stack costs nothing to hold: $3{,}584$
bytes per block against $57{,}376$ at $1024$, and the kernel is register-limited
to two blocks per multiprocessor at either capacity, so what the reduction buys
is balance and not occupancy.

\emph{Bound refresh interval.} The running minimum is re-read from global memory
on \emph{every} iteration of the device loop. Refreshing every $4$th, $16$th,
$64$th or $256$th iteration instead was monotonically worse on all three scenes
tested --- on cloth-funnel, $638.7\,\mathrm{ms}$ at every iteration against
$3377.5\,\mathrm{ms}$ at every $256$th --- because the pruning a fresh bound buys
exceeds what the atomic costs.

\emph{Unit of work per query.} One thread per query beat one block per query by
between $2.76\times$ and $4.47\times$ on the three scenes, with identical false
positives and no missed collision, because the block-per-query kernel was bound
by \emph{scheduling blocks} and not by the search: cloth-funnel issues $843{,}414$
of them, each classifying $63$ boxes across $128$ threads.

\emph{Work-list capacity.} This is the one of the five that shows up as a loss
of accuracy rather than as a slowdown, and it is worth separating because the
symptom misleads. A capacity sized from the mean is far too small for the tail. The host explores about $1.2$ boxes
per query when a shared bound is pruning and about $11$ when it is not, which
suggests a stack of a few dozen entries. On a vertex-quad query whose exact
answer is $2/3$, a $32$-entry stack reports $0.6655369599$ where $128$ and $512$
entries both report $0.6666666238$: the answer falls $1.13\times10^{-3}$ before
the true value where the triangle path reaches $3.0\times10^{-7}$. All three
values are early, so the guarantee holds throughout and only the step a solver
can take is lost --- but read as a conservativeness failure it sends an
implementer to the acceptance test, where the fault is not.

\subsection{Why the work divides the way it does}
\label{sec:parallel:why}

\Cref{tab:difficulty} counts, for every candidate pair of one scene, the boxes
each narrow phase examined before it answered. Both distributions are
heavy-tailed, which is the premise both kernels were built on: seven queries in
ten are settled by one box or none, while a few need tens of thousands and the
worst needs millions. The two differ in where the weight sits, and the
difference runs in opposite directions at the two ends of the distribution.

In total the device examines $4.8\times$ as many boxes as the host,
$106.9$ million against $22.2$, and the excess is in the \emph{middle} of the
distribution rather than in the tail: $1.33$ million device queries need at
least eight boxes against $0.82$ million host queries, and the ratio is similar
at $64$ and at $1{,}024$. A per-box factor of this kind is what
\cref{alg:host,alg:device} predict, since the device re-evaluates all eight
corners of both children at every split where the host inherits four of each
child's eight from its parent (\cref{sec:parallel:host}). That factor falls on
the operation which dominates both kernels and is paid on every box.

At the far end of the distribution the ordering reverses, and it is the host that
owns the tail. Six host queries examine more than a million boxes and its worst
examines $4{,}591{,}271$; no device query reaches a million and its worst
examines $101{,}246$, $45\times$ fewer. This is the bound schedule rather than
the search. A host block is seeded with the shared minimum as it stands when that
block starts, so a block the scheduler runs early searches a window in $t$ that
no other block has yet narrowed, and a hard query landing in such a block is
searched almost unpruned. The device issues every query of the batch in one
launch, so the minimum collapses while the hard queries are still being searched,
and spilling bounds how long any one block can be held by a single query.

Neither effect is a property of the hardware, and the second of them is the
reason the device wins the narrow phase on the scenes with the most candidate
pairs while losing it on the smallest (\cref{sec:results:proc}): the collapse of
the shared bound within a launch is worth more the more queries the launch
carries.

The broad phase inverts both arguments, having no sequential dependence and no
shared bound to lose. The GPU accordingly wins that phase on all six scenes, and
the pipeline wins end to end on four of them despite losing the narrow phase on
four. \Cref{sec:results} measures all of it.

\begin{table}[htbp]
  \centering
  \caption{Boxes examined per query on cloth-funnel, over every one of the
    $25{,}192{,}698$ candidate pairs the broad phase produced in the
    scene's $577$ steps. Both processors are handed the same
    candidates and asked for one earliest time of impact per step, so both
    prune against a shared running minimum. Each narrow phase counts its own
    boxes per query and buckets them by $\lfloor\log_{2}\rfloor$ with the same
    code, so the two columns measure the same quantity; the device's count also
    includes boxes it evaluated and then discarded because the bound had already
    passed them, which the host discards before counting, and which is
    $0.06\%$ of the device total here. Timings are not reported from this
    run: the counters serialise the device kernel.}
  \label{tab:difficulty}
  \begin{tabular}{lrr}
    \toprule
    boxes examined & host queries & device queries \\
    \midrule
    $0$--$1$                 & $18{,}852{,}925$ & $17{,}220{,}209$ \\
    $\ge 8$                  & $817{,}187$ & $1{,}332{,}867$ \\
    $\ge 64$                 & $23{,}035$ & $43{,}243$ \\
    $\ge 1{,}024$            & $2{,}046$ & $6{,}383$ \\
    $\ge 16{,}384$           & $92$ & $61$ \\
    $\ge 1{,}048{,}576$      & $6$ & $0$ \\
    \midrule
    worst single query       & $4{,}591{,}271$ & $101{,}246$ \\
    \midrule
    boxes in total           & $22{,}195{,}731$ & $106{,}911{,}080$ \\
    mean per query           & $0.88$ & $4.24$ \\
    \bottomrule
  \end{tabular}
  \par\smallskip\footnotesize Source: \texttt{benchmark/results/boxes/}
\end{table}

